\subsection{3.1 主路通行能力}

主路通行能力占总分 25\%，重点是“前方遇静止目标物”，二级权重 50\%，折算到总分就是 12.5\%。这说明 C-ICAP 对城市 NOA 的第一要求是：系统必须能识别并处理城市道路上突然出现在前方的障碍物。

典型场景包括

这里的评价不是只看有没有撞上。很多场景还会看安全、舒适、效率三件事。比如系统换道避让障碍物，可以拿到效率分；如果只是在本车道刹停，安全分可能有，但效率分会低。最大纵向减速度、最大横向加速度也会影响舒适分。一般来说，纵向减速度不超过 2 m/s²、横向加速度不超过 2 m/s² 是比较舒适的；超过 4 m/s² 的急减速或超过 3 m/s² 的横向加速度，分数会明显下降甚至为 0。

以城市领航里常见的“前方遇静止目标物”为例，通常是：

\texttt{单项得分 = 安全 × 50\% + 舒适 × 25\% + 效率 × 25\%}
